\documentclass{article}
\usepackage[all]{xy}
\usepackage{amsmath,amsthm,amssymb,color,latexsym}
\usepackage{geometry}        
\geometry{letterpaper}    
\usepackage{graphicx}
\usepackage[]{inputenc}
\usepackage[T1]{fontenc}
\usepackage{fullpage}
\usepackage{coqdoc}
\usepackage{amsmath,amssymb}
\usepackage{url}

\begin{document}
	
	\section*{Charles Pinter - A Book of Set Theory}
	\subsection*{Chapter 1}
	\par\bigskip
	\subsection*{Exercises 1.1}
	\begin{enumerate}
		\item \begin{proof}
			For all sentences $P,Q$ and $R$, we are going to proof the following statments to be true, using its associated truth table. First of all, notice the truth table for the expression $P \iff Q $ (that is, the abbreviaton of $(P \Rightarrow Q) \land (P \Rightarrow P)$): 
			
			\begin{center}
				\begin{tabular}{|c|c|c|c|c|}
					\hline
					$P$ & $Q$ & $P \Rightarrow Q$ & $Q \Rightarrow P$ & $P \iff Q: \tiny (P \Rightarrow Q) \land (P \Rightarrow P)$ \\
					\hline
					1 & 1 & 1 & 1 & 1\\
					1 & 0 & 1 & 0 & 0\\
					0 & 1 & 0 & 1 & 0\\
					0 & 0 & 1 & 1 & 1\\
					\hline
				\end{tabular}
			\end{center} 
			
			{
				\begin{enumerate}
					\item[i.] $P \lor Q \iff Q \lor P$
					
					\begin{center}
						\begin{tabular}{|c|c|c|c|c|}
							\hline
							$P$ & $Q$ & $P \lor Q$ & $Q \lor P$ & $P \lor Q \iff Q \lor P $ \\
							\hline
							1 & 1 & 1 & 1 & 1\\
							1 & 0 & 1 & 1 & 1\\
							0 & 1 & 1 & 1 & 1\\
							0 & 0 & 0 & 0 & 1\\
							\hline
						\end{tabular}
					\end{center}
					
					\item[ii.] $P \lor (Q \lor R) \iff (P \lor Q) \lor R$
					
					\begin{center}
						\begin{tabular}{|c|c|c|c|c|c|c|c|}
							\hline
							$P$ & $Q$ & $R$ & $Q \lor R$ & $P \lor Q$ & $P \lor (Q \lor R)$ & $(P \lor Q) \lor R$ & $P \lor (Q \lor R) \iff (P \lor Q) \lor R$ \\
							\hline
							1 & 1 & 1 & 1 & 1 & 1 & 1 & 1 \\
							1 & 1 & 0 & 1 & 1 & 1 & 1 & 1 \\
							1 & 0 & 1 & 1 & 1 & 1 & 1 & 1 \\
							1 & 0 & 0 & 0 & 1 & 1 & 1 & 1 \\
							0 & 1 & 1 & 1 & 1 & 1 & 1 & 1 \\
							0 & 1 & 0 & 1 & 1 & 1 & 1 & 1 \\
							0 & 0 & 1 & 1 & 0 & 1 & 1 & 1 \\
							0 & 0 & 0 & 0 & 0 & 0 & 0 & 1 \\
							\hline
						\end{tabular}
					\end{center} 
					
					\item[iii.]$P \land (Q \lor R) \iff (P \land Q) \lor (P \land R)$.
					
					Let $X = P \land (Q \lor R)$ and $Y = (P \land Q) \lor (P \land R)$.
					
					\begin{center}
						\begin{tabular}{|c|c|c|c|c|c|c|c|c|c|}
							\hline
							$P$ & $Q$ & $R$ & $Q \lor R$ & $P \lor Q$ & $P \land (Q \lor R)$ & $P \land Q$ & $P \land R$ & $(P \land Q) \lor (P \land R)$ & $X \iff Y$ \\
							\hline
							1 & 1 & 1 & 1 & 1 & 1 & 1 & 1 & 1 & 1 \\
							1 & 1 & 0 & 1 & 1 & 1 & 1 & 0 & 1& 1 \\
							1 & 0 & 1 & 1 & 1 & 1 & 0 & 1 & 1 & 1 \\
							1 & 0 & 0 & 0 & 0 & 0 & 0 & 0 & 0 & 1\\
							0 & 1 & 1 & 1 & 0 & 0 & 0 & 0 & 0 &  1 \\
							0 & 1 & 0 & 1 & 0 & 0 & 0 & 0 & 0 & 1\\
							0 & 0 & 1 & 1 & 0 & 0 & 0 & 0 & 0 & 1\\
							0 & 0 & 0 & 0 & 0 & 0 & 0 & 0 & 0 &  1\\
							\hline
						\end{tabular}
					\end{center} 
					
					\item[iv.] $P \lor P \iff P$ 
					
					\begin{center}
						\begin{tabular}{|c|c|c|c|}
							\hline
							$P$ & $P$ & $P \lor P$ & $P \lor P \iff P$ \\
							\hline
							1 & 1 & 1 & 1 \\
							O & 0 & 0 & 1 \\
							\hline
						\end{tabular}
					\end{center}
				\end{enumerate}
			}  
		\end{proof}
		\item \begin{proof}
			Let's proof DeMorgan's laws using its derive truth table:
			
			\begin{enumerate}
				\item[a)] $\neg (P  \lor Q) \iff \neg P \land \neg Q$
				
				\begin{center}
					\begin{tabular}{|c|c|c|c|c|c|c|c|}
						\hline
						$P$ & $Q$ & $P \lor Q$ & $\neg P$ & $\neg Q$ & $\neg(P \lor Q)$ & $\neg P \land \neg Q$ & $  \neg(P \lor Q) \iff \neg P \land \neg Q$ \\
						\hline
						1 & 1 & 1 & 0 & 0 & 0 & 0 & 1 \\
						1 & 0 & 1 & 0 & 1 & 0 & 0 & 1 \\
						0 & 1 & 1 & 1 & 0 & 0 & 0 & 1 \\
						0 & 0 & 0 & 1 & 1 & 1 & 1 & 1 \\
						\hline
					\end{tabular}
				\end{center}
				
				\item[b)] $\neg (P  \land Q) \iff \neg P \lor \neg Q$
				
				\begin{center}
					\begin{tabular}{|c|c|c|c|c|c|c|c|}
						\hline
						$P$ & $Q$ & $P \land Q$ & $\neg P$ & $\neg Q$ & $\neg(P \land Q)$ & $\neg P \lor \neg Q$ & $  \neg(P \land Q) \iff \neg P \lor \neg Q$ \\
						\hline
						1 & 1 & 1 & 0 & 0 & 0 & 0 & 1 \\
						1 & 0 & 0 & 0 & 1 & 1 & 1 & 1 \\
						0 & 1 & 0 & 1 & 0 & 1 & 1 & 1 \\
						0 & 0 & 0 & 1 & 1 & 1 & 1 & 1 \\
						\hline
					\end{tabular}
				\end{center} 
			\end{enumerate}
		\end{proof}
		
		\item \begin{proof}
			\begin{enumerate}
				\item[a)] $\neg \neg P \Rightarrow P$
				
				\begin{center}
					\begin{tabular}{|c|c|c|c|}
						\hline
						$P$ & $\neg P$ & $\neg \neg P$ & $\neg \neg P \Rightarrow P$ \\
						\hline
						1 & 0 & 1 & 1 \\
						O & 1 & 0 & 1 \\
						\hline
					\end{tabular}
					\end{center}
				\item[b)] $P \Rightarrow \neg \neg P$ 
				
				\begin{center}
					\begin{tabular}{|c|c|c|c|}
						\hline
						$P$ & $\neg P$ & $\neg \neg P$ & $ P \Rightarrow \neg \neg P$ \\
						\hline
						1 & 0 & 1 & 1 \\
						O & 1 & 0 & 1 \\
						\hline
					\end{tabular}
				\end{center}
				
			\end{enumerate}
		\end{proof}
		
		\item \begin{proof} 
			
			\begin{enumerate}
			\item[a)] $(P \Rightarrow Q) \iff \neg Q \Rightarrow \neg P$
			
			\begin{center}
				\begin{tabular}{|c|c|c|c|c|c|c|}
					\hline
					$P$ & $Q$ & $\neg Q$ & $\neg P$ & $P \Rightarrow Q$ & $\neg Q \Rightarrow \neg P$ & $ (P \Rightarrow Q) \iff \neg (Q \Rightarrow \neg P)$ \\
					\hline
					1 & 1 & 0 & 0 & 1 & 1 & 1  \\
					1 & 0 & 1 & 0 & 0 & 0 & 1  \\
					0 & 1 & 0 & 1 & 1 & 1 & 1  \\
					0 & 0 & 1 & 1 & 1 & 1 & 1  \\
					\hline
				\end{tabular}
			\end{center}
			
			\par\bigskip
			
			\item[b)] $(P \Rightarrow Q) \iff (\neg P \lor Q)$
			
			\begin{center}
				\begin{tabular}{|c|c|c|c|c|c|}
					\hline
					$P$ & $Q$ & $\neg P$ & $P \Rightarrow Q$ & $\neg P \lor Q$ & $ (P \Rightarrow Q) \iff \neg (Q \Rightarrow \neg P)$ \\
					\hline
					1 & 1 & 0 & 1 & 1 & 1   \\
					1 & 0 & 0 & 0 & 0 & 1   \\
					0 & 1 & 1 & 1 & 1 & 1   \\
					0 & 0 & 1 & 1 & 1 & 1   \\
					\hline
				\end{tabular}
			\end{center}
			
			\item[c)] $(P \Rightarrow Q) \iff \neg(P \land \neg Q)$
			
			\begin{center}
				\begin{tabular}{|c|c|c|c|c|c|c|}
					\hline
					$P$ & $Q$ & $\neg Q$ & $P \Rightarrow Q$ & $P \land \neg Q$ & $\neg (P \land \neg Q)$& $ (P \Rightarrow Q) \iff \neg (P \land \neg Q)$ \\
					\hline
					1 & 1 & 0 & 1 & 0 & 1  & 1 \\
					1 & 0 & 1 & 0 & 1 & 0  & 1\\
					0 & 1 & 0 & 1 & 0 & 1  & 1\\
					0 & 0 & 1 & 1 & 0 & 1  & 1\\
					\hline
				\end{tabular}
			\end{center}
			
			\item[d)] $(P \land (P \Rightarrow Q)) \Rightarrow Q$
			
			\begin{center}
				\begin{tabular}{|c|c|c|c|c|}
					\hline
					$P$ & $Q$ & $P \Rightarrow Q$ & $P \land (P \Rightarrow Q)$ & $(P \land (P \Rightarrow Q)) \Rightarrow Q$ \\
					\hline
					1 & 1 & 1 & 1 & 1  \\
					1 & 0 & 0 & 0 & 1 \\
					0 & 1 & 1 & 0 & 1 \\
					0 & 0 & 1 & 0 & 1 \\
					\hline
				\end{tabular}
			\end{center}
			
			\item[e)] $((P \lor Q) \land \neg P) \Rightarrow Q$
			
			\begin{center}
				\begin{tabular}{|c|c|c|c|c|c|}
					\hline
					$P$ & $Q$ & $\neg P$ & $P \lor Q$ & $((P \lor Q) \land \neg P)$ & $((P \lor Q) \land \neg P) \Rightarrow Q$\\
					\hline
					1 & 1 & 1 & 1 & 1 & 1\\
					1 & 0 & 0 & 0 & 0 & 1\\
					0 & 1 & 1 & 0 & 0 & 1\\
					0 & 0 & 1 & 0 & 0 & 1\\
					\hline
				\end{tabular}
			\end{center}
		\end{enumerate}
		\end{proof}
		\item \begin{proof}
			\begin{enumerate}
				\item[a)] $((P \Rightarrow Q) \land (Q \Rightarrow R)) \Rightarrow (P \Rightarrow R)$
				
				Let $X = (P \Rightarrow Q) \land (Q \Rightarrow R)$
				
				
				\begin{center}
					\begin{tabular}{|c|c|c|c|c|c|c|c|}
						\hline
						$P$ & $Q$ & $R$ & $P \Rightarrow Q$ & $Q \Rightarrow R$ & $P \Rightarrow R$ & $(P \Rightarrow Q) \land (Q \Rightarrow R)$ & $X \Rightarrow (P \Rightarrow R)$ \\
						\hline
						1 & 1 & 1 & 1 & 1 & 1 & 1 & 1 \\
						1 & 1 & 0 & 1 & 0 & 0 & 0 & 1 \\
						1 & 0 & 1 & 0 & 1 & 1 & 0 & 1 \\
						1 & 0 & 0 & 0 & 1 & 0 & 0 & 1 \\
						0 & 1 & 1 & 1 & 1 & 1 & 1 & 1 \\
						0 & 1 & 0 & 1 & 0 & 1 & 0 & 1 \\
						0 & 0 & 1 & 1 & 1 & 1 & 0 & 1 \\
						0 & 0 & 0 & 1 & 1 & 1 & 0 & 1 \\
						\hline
					\end{tabular}
				\end{center}
				
				\item[b)] $((P \Rightarrow Q) \land (R \Rightarrow Q)) \iff ((P \lor R) \Rightarrow Q)$
				
				Let $X = ((P \Rightarrow Q) \land (R \Rightarrow Q))$
				
				
				\begin{center}
					\begin{tabular}{|c|c|c|c|c|c|c|c|c|}
						\hline
						$P$ & $Q$ & $R$ & $P \Rightarrow Q$ & $R \Rightarrow Q$ & $X$ & $P \lor R$ & $((P \lor R) \Rightarrow Q)$ & $X \iff ((P \lor R) \Rightarrow Q)$ \\
						\hline
						1 & 1 & 1 & 1 & 1 & 1 & 1 & 1 & 1\\
						1 & 1 & 0 & 1 & 1 & 1 & 1 & 1 & 1\\
						1 & 0 & 1 & 0 & 0 & 0 & 1 & 0 & 1\\
						1 & 0 & 0 & 0 & 1 & 0 & 1 & 0 & 1\\
						0 & 1 & 1 & 1 & 1 & 1 & 1 & 1 & 1\\
						0 & 1 & 0 & 1 & 1 & 1 & 0 & 1 & 1\\
						0 & 0 & 1 & 1 & 0 & 0 & 1 & 0 & 1\\
						0 & 0 & 0 & 1 & 1 & 1 & 0 & 1 & 1\\
						\hline
					\end{tabular}
				\end{center}
				
				
				\item[c)] $((P \Rightarrow Q) \land (P \Rightarrow R)) \iff (P \Rightarrow (Q \land R))$
				
				Let $X = ((P \Rightarrow Q) \land (P \Rightarrow R)$
				
				\begin{center}
					\begin{tabular}{|c|c|c|c|c|c|c|c|c|}
						\hline
						$P$ & $Q$ & $R$ & $P \Rightarrow Q$ & $P \Rightarrow R$ & $X$ & $Q \land R$ & $P \Rightarrow (Q \land R)$ & $X \iff (P \Rightarrow (Q \land R))$ \\
						\hline
						1 & 1 & 1 & 1 & 1 & 1 & 1 & 1 & 1\\
						1 & 1 & 0 & 1 & 0 & 0 & 0 & 0 & 1\\
						1 & 0 & 1 & 0 & 1 & 0 & 0 & 0 & 1\\
						1 & 0 & 0 & 0 & 0 & 0 & 0 & 0 & 1\\
						0 & 1 & 1 & 1 & 1 & 1 & 1 & 1 & 1\\
						0 & 1 & 0 & 1 & 1 & 1 & 0 & 1 & 1\\
						0 & 0 & 1 & 1 & 1 & 1 & 0 & 1 & 1\\
						0 & 0 & 0 & 1 & 1 & 1 & 0 & 1 & 1\\
						\hline
					\end{tabular}
				\end{center}
			\end{enumerate}
		\end{proof}
		\item \begin{proof}

			\begin{enumerate}
				\item[a)] $P \lor Q \iff P \lor R$
				
				Supoose that for all sentences $Q$, $R$ and $P$, $Q \iff R$ is true. Since $Q \iff R$ is the abbreviation of $(Q \Rightarrow R) \land (R \Rightarrow Q)$, we can assume that $R \Rightarrow Q$ is true. Then, let's derive the truth table for $P \lor R \Rightarrow P \lor Q$
				
					\begin{center}
					\begin{tabular}{|c|c|c|c|c|c|}
						\hline
						$P$ & $Q$ & $R$ & $P \lor R$ & $P \lor Q$ & $((P \lor R) \Rightarrow (P \lor Q)$\\
						\hline
						1 & 1 & 1 & 1 & 1 & 1\\
						1 & 1 & 0 & 1 & 1 & 1\\
						1 & 0 & 1 & 1 & 1 & 1\\
						1 & 0 & 0 & 1 & 1 & 1\\
						0 & 1 & 1 & 1 & 1 & 1\\
						0 & 1 & 0 & 0 & 1 & 1\\
						0 & 0 & 1 & 1 & 0 & 0\\
						0 & 0 & 0 & 0 & 0 & 1\\
						\hline
					\end{tabular}
				\end{center}
				
				In this maner, since $R \Rightarrow Q$ is true, we cannot have simultaneuously, R true and Q false, then we may disregard the seventh line of the table. In all the remaining lines, $P \lor Q \Rightarrow P \lor R$ takes the value 1 (true). Finally, because $Q \Rightarrow R$ is also true, Theorem 1.7 also holds, then $P \lor Q \Rightarrow Q \lor R$ and $P \lor R \Rightarrow Q \lor R$ are true, that is $P \lor Q \iff P \lor R$ is true.
				
				\item[b)] Analogous to the above. 
				\item[c)] $(P \Rightarrow Q) \iff (P \Rightarrow R)$.
				
				Suppose $Q \iff R$ is true. Let derive the truth table for $(P \Rightarrow Q) \iff  (P \Rightarrow R)$
				
				\begin{center}
					\begin{tabular}{|c|c|c|c|c|c|}
						\hline
						$P$ & $Q$ & $R$ & $P \Rightarrow Q$ & $P \Rightarrow R$ & $(P \Rightarrow R) \iff (P \Rightarrow Q)$\\
						\hline
						1 & 1 & 1 & 1 & 1 & 1\\
						1 & 1 & 0 & 0 & 0 & 0\\
						1 & 0 & 1 & 0 & 1 & 0\\
						1 & 0 & 0 & 1 & 0 & 1\\
						0 & 1 & 1 & 1 & 1 & 1\\
						0 & 1 & 0 & 1 & 1 & 1\\
						0 & 0 & 1 & 1 & 1 & 1\\
						0 & 0 & 0 & 1 & 1 & 1\\
						\hline
					\end{tabular}
				\end{center}
				Since $Q \Rightarrow R$ and $R \Rightarrow Q$ is true, we cannot have simultaneously $Q$ true and $R$ false, and $P$ true and $R$ false. Then, we may disregard the second and third line of the table. Finally, in all the remaining lines $(P \Rightarrow Q) \iff (P \Rightarrow R)$ takes the value 1 (true).
			\end{enumerate}
		\end{proof}
		
		\item \begin{proof} Suppose $P \Rightarrow Q$ and $R \Rightarrow S$ are true. 
			
			\begin{enumerate}	
			  \item[a)] $P \lor R \Rightarrow Q \lor S$ is true.
			  
			  Since $P \Rightarrow Q$ is true, for Theorem 1.7 we have that 
			  
			  \begin{equation*}
			  	P \lor R \iff R \lor P \Rightarrow R \lor Q \iff Q \lor R
			  \end{equation*}
			  
			  and, since $R \Rightarrow S$ is true, then $Q \lor R \Rightarrow Q \lor S$ is true. Therefore, for exercise 5. (a) we have that 
			  
			  \begin{equation*}
			  	((P \lor R \Rightarrow Q \lor R) \land (Q \lor R \Rightarrow Q \lor S)) \Rightarrow (P \lor R \Rightarrow Q \lor S)
			  \end{equation*}
			  
			  is true for all $P$, $Q$, $R$ and $S$. Then, we cannot have simultaneously $((P \lor R \Rightarrow Q \lor R) \land (Q \lor R \Rightarrow Q \lor S))$ true and $(P \lor R \Rightarrow Q \lor S)$ false. But $((P \lor R \Rightarrow Q \lor R) \land (Q \lor R \Rightarrow Q \lor S))$ is always true, therefore $(P \lor R \Rightarrow Q \lor S)$ has to be true.
			  
			  \item[b)] Analogous to the prevoius one, using Theorem 1.7 (ii).
			\end{enumerate}
			
		\end{proof}
		
		
			
	\end{enumerate}
	
  \section*{Charles Pinter - A Book of Set Theory (Formal Methods)}
	\subsection*{Chapter 1}
	\par\bigskip
	\subsection*{Exercises 1.1}
  \coqlibrary{chapter 1 exercises 1 point 1}{Library }{chapter\_1\_exercises\_1\_point\_1}

\begin{coqdoccode}
\end{coqdoccode}
This is basically checking for every possibility. It simply splits the directions and the options and then applies them one by one.

\begin{coqdoccode}
\coqdocnoindent
\coqdockw{Theorem} \coqdocvar{theorem\_1\_8\_i} : \coqdockw{\ensuremath{\forall}} (\coqdocvar{P} \coqdocvar{Q}: \coqdockw{Prop}), \coqdocvar{P} \ensuremath{\lor} \coqdocvar{Q} \ensuremath{\leftrightarrow} \coqdocvar{Q} \ensuremath{\lor} \coqdocvar{P}.\coqdoceol
\coqdocnoindent
\coqdockw{Proof}.\coqdoceol
\coqdocindent{1.00em}
\coqdoctac{intros} \coqdocvar{P} \coqdocvar{Q}. \coqdoctac{split}; \coqdoctac{intros} [\coqdocvar{H1} \ensuremath{|} \coqdocvar{H2}].\coqdoceol
\coqdocindent{2.00em}
** \coqdoctac{right}. \coqdoctac{apply} \coqdocvar{H1}.\coqdoceol
\coqdocindent{2.00em}
** \coqdoctac{left}. \coqdoctac{apply} \coqdocvar{H2}.\coqdoceol
\coqdocindent{2.00em}
** \coqdoctac{right}. \coqdoctac{apply} \coqdocvar{H1}.\coqdoceol
\coqdocindent{2.00em}
** \coqdoctac{left}. \coqdoctac{apply} \coqdocvar{H2}.\coqdoceol
\coqdocnoindent
\coqdockw{Qed}.\coqdoceol
\coqdocemptyline
\end{coqdoccode}
Similar to the last point, it splits the possibilities and applies them as needed.

\begin{coqdoccode}
\coqdocemptyline
\coqdocnoindent
\coqdockw{Theorem} \coqdocvar{theorem\_1\_8\_ii} : \coqdockw{\ensuremath{\forall}} (\coqdocvar{P} \coqdocvar{Q} \coqdocvar{R} : \coqdockw{Prop}),\coqdoceol
\coqdocindent{1.00em}
\coqdocvar{P} \ensuremath{\lor}  (\coqdocvar{Q} \ensuremath{\lor} \coqdocvar{R}) \ensuremath{\leftrightarrow} (\coqdocvar{P} \ensuremath{\lor} \coqdocvar{Q}) \ensuremath{\lor} \coqdocvar{R}.\coqdoceol
\coqdocnoindent
\coqdockw{Proof}.\coqdoceol
\coqdocindent{1.00em}
\coqdoctac{intros} \coqdocvar{P} \coqdocvar{Q} \coqdocvar{R}. \coqdoctac{split}.\coqdoceol
\coqdocindent{1.00em}
\ensuremath{\times} \coqdoctac{intros} [\coqdocvar{HP} \ensuremath{|} [\coqdocvar{HQ} \ensuremath{|} \coqdocvar{HR}]].\coqdoceol
\coqdocindent{2.00em}
** \coqdoctac{left}. \coqdoctac{left}. \coqdoctac{apply} \coqdocvar{HP}.\coqdoceol
\coqdocindent{2.00em}
** \coqdoctac{left}. \coqdoctac{right}. \coqdoctac{apply} \coqdocvar{HQ}.\coqdoceol
\coqdocindent{2.00em}
** \coqdoctac{right}. \coqdoctac{apply} \coqdocvar{HR}.\coqdoceol
\coqdocindent{1.00em}
\ensuremath{\times} \coqdoctac{intros} [[\coqdocvar{HP} \ensuremath{|} \coqdocvar{HQ}] \ensuremath{|} \coqdocvar{HR}].\coqdoceol
\coqdocindent{2.00em}
** \coqdoctac{left}. \coqdoctac{apply} \coqdocvar{HP}.\coqdoceol
\coqdocindent{2.00em}
** \coqdoctac{right}. \coqdoctac{left}. \coqdoctac{apply} \coqdocvar{HQ}.\coqdoceol
\coqdocindent{2.00em}
** \coqdoctac{right}. \coqdoctac{right}. \coqdoctac{apply} \coqdocvar{HR}.\coqdoceol
\coqdocnoindent
\coqdockw{Qed}.\coqdoceol
\coqdocemptyline
\end{coqdoccode}
You will become accustomed to the way it works. Simply splitting the possibilities and applying them is basically what is done here.

\begin{coqdoccode}
\coqdocemptyline
\coqdocnoindent
\coqdockw{Theorem} \coqdocvar{theorem\_1\_8\_iii} : \coqdockw{\ensuremath{\forall}} (\coqdocvar{P} \coqdocvar{Q} \coqdocvar{R} : \coqdockw{Prop}), \coqdocvar{P} \ensuremath{\land}  (\coqdocvar{Q} \ensuremath{\lor} \coqdocvar{R}) \ensuremath{\leftrightarrow} (\coqdocvar{P} \ensuremath{\land} \coqdocvar{Q}) \ensuremath{\lor} (\coqdocvar{P} \ensuremath{\land} \coqdocvar{R}).\coqdoceol
\coqdocnoindent
\coqdockw{Proof}.\coqdoceol
\coqdocindent{1.00em}
\coqdoctac{intros} \coqdocvar{P} \coqdocvar{Q} \coqdocvar{R}. \coqdoctac{split}.\coqdoceol
\coqdocindent{1.00em}
\ensuremath{\times} \coqdoctac{intros} [\coqdocvar{HP} [\coqdocvar{HQ} \ensuremath{|} \coqdocvar{HR}]].\coqdoceol
\coqdocindent{2.00em}
** \coqdoctac{left}. \coqdoctac{split}. \coqdoctac{apply} \coqdocvar{HP}. \coqdoctac{apply} \coqdocvar{HQ}.\coqdoceol
\coqdocindent{2.00em}
** \coqdoctac{right}. \coqdoctac{split}. \coqdoctac{apply} \coqdocvar{HP}. \coqdoctac{apply} \coqdocvar{HR}.\coqdoceol
\coqdocindent{1.00em}
\ensuremath{\times} \coqdoctac{intros}[[\coqdocvar{HP} \coqdocvar{HQ}] \ensuremath{|} [\coqdocvar{HP} \coqdocvar{HR}]].\coqdoceol
\coqdocindent{2.00em}
** \coqdoctac{split}. \coqdoctac{apply} \coqdocvar{HP}. \coqdoctac{left}. \coqdoctac{apply} \coqdocvar{HQ}.\coqdoceol
\coqdocindent{2.00em}
** \coqdoctac{split}. \coqdoctac{apply} \coqdocvar{HP}. \coqdoctac{right}. \coqdoctac{apply} \coqdocvar{HR}.\coqdoceol
\coqdocnoindent
\coqdockw{Qed}.\coqdoceol
\coqdocemptyline
\end{coqdoccode}
One important note in this exercise (one that I strongly recommend looking into) is the use of ; in this proof. When you use it after opening different branches, it applies the command in both branches (really cool for solving many problems at the same time).

\begin{coqdoccode}
\coqdocemptyline
\coqdocnoindent
\coqdockw{Theorem} \coqdocvar{theorem\_1\_8\_iv} : \coqdockw{\ensuremath{\forall}} (\coqdocvar{P} : \coqdockw{Prop}), \coqdocvar{P} \ensuremath{\lor} \coqdocvar{P} \ensuremath{\leftrightarrow} \coqdocvar{P}.\coqdoceol
\coqdocnoindent
\coqdockw{Proof}.\coqdoceol
\coqdocindent{1.00em}
\coqdoctac{intros} \coqdocvar{P}. \coqdoctac{split}.\coqdoceol
\coqdocindent{1.00em}
\ensuremath{\times} \coqdoctac{intros} [\coqdocvar{HP} \ensuremath{|} \coqdocvar{HP}]; \coqdoctac{apply} \coqdocvar{HP}.\coqdoceol
\coqdocindent{1.00em}
\ensuremath{\times} \coqdoctac{intros} \coqdocvar{HP}; \coqdoctac{left}; \coqdoctac{apply} \coqdocvar{HP}.\coqdoceol
\coqdocnoindent
\coqdockw{Qed}.\coqdoceol
\coqdocemptyline
\end{coqdoccode}
As before, it is simply applying every option. Note that it is simply reversing the order of application.

\begin{coqdoccode}
\coqdocemptyline
\coqdocnoindent
\coqdockw{Theorem} \coqdocvar{theorem\_1\_8\_i'} : \coqdockw{\ensuremath{\forall}} (\coqdocvar{P} \coqdocvar{Q}: \coqdockw{Prop}), \coqdocvar{P} \ensuremath{\land} \coqdocvar{Q} \ensuremath{\leftrightarrow} \coqdocvar{Q} \ensuremath{\land} \coqdocvar{P}.\coqdoceol
\coqdocnoindent
\coqdockw{Proof}.\coqdoceol
\coqdocindent{1.00em}
\coqdoctac{intros} \coqdocvar{P} \coqdocvar{Q}. \coqdoctac{split}.\coqdoceol
\coqdocindent{1.00em}
\ensuremath{\times} \coqdoctac{intros} [\coqdocvar{HP} \coqdocvar{HQ}]. \coqdoctac{split}. \coqdoctac{apply} \coqdocvar{HQ}. \coqdoctac{apply} \coqdocvar{HP}.\coqdoceol
\coqdocindent{1.00em}
\ensuremath{\times} \coqdoctac{intros} [\coqdocvar{HQ} \coqdocvar{HP}]. \coqdoctac{split}. \coqdoctac{apply} \coqdocvar{HP}. \coqdoctac{apply} \coqdocvar{HQ}.\coqdoceol
\coqdocnoindent
\coqdockw{Qed}.\coqdoceol
\coqdocemptyline
\end{coqdoccode}
You can (and I actually did) try to simplify this proof with ; and try, however it is uglier and more convoluted that way.

\begin{coqdoccode}
\coqdocemptyline
\coqdocnoindent
\coqdockw{Theorem} \coqdocvar{theorem\_1\_8\_ii'} : \coqdockw{\ensuremath{\forall}} (\coqdocvar{P} \coqdocvar{Q} \coqdocvar{R} : \coqdockw{Prop}), \coqdocvar{P} \ensuremath{\land}  (\coqdocvar{Q} \ensuremath{\land} \coqdocvar{R}) \ensuremath{\leftrightarrow} (\coqdocvar{P} \ensuremath{\land} \coqdocvar{Q}) \ensuremath{\land} \coqdocvar{R}.\coqdoceol
\coqdocnoindent
\coqdockw{Proof}.\coqdoceol
\coqdocindent{1.00em}
\coqdoctac{intros} \coqdocvar{P} \coqdocvar{Q} \coqdocvar{R}. \coqdoctac{split}.\coqdoceol
\coqdocindent{1.00em}
\ensuremath{\times} \coqdoctac{intros} [\coqdocvar{HP} [\coqdocvar{HQ} \coqdocvar{HR}]]. \coqdoctac{split}. \coqdoctac{split}. \coqdoctac{apply} \coqdocvar{HP}. \coqdoctac{apply} \coqdocvar{HQ}. \coqdoctac{apply} \coqdocvar{HR}.\coqdoceol
\coqdocindent{1.00em}
\ensuremath{\times} \coqdoctac{intros} [[\coqdocvar{HP} \coqdocvar{HQ}] \coqdocvar{HR}]. \coqdoctac{split}. \coqdoctac{apply} \coqdocvar{HP}. \coqdoctac{split}. \coqdoctac{apply} \coqdocvar{HQ}. \coqdoctac{apply} \coqdocvar{HR}.\coqdoceol
\coqdocnoindent
\coqdockw{Qed}.\coqdoceol
\coqdocemptyline
\end{coqdoccode}
It is in the same direction as every other proof in this first exercise. Simply working every possibility.

\begin{coqdoccode}
\coqdocemptyline
\coqdocnoindent
\coqdockw{Theorem} \coqdocvar{theorem\_1\_8\_iii'} : \coqdockw{\ensuremath{\forall}} (\coqdocvar{P} \coqdocvar{Q} \coqdocvar{R} : \coqdockw{Prop}), \coqdocvar{P} \ensuremath{\lor}  (\coqdocvar{Q} \ensuremath{\land} \coqdocvar{R}) \ensuremath{\leftrightarrow} (\coqdocvar{P} \ensuremath{\lor} \coqdocvar{Q}) \ensuremath{\land} (\coqdocvar{P} \ensuremath{\lor} \coqdocvar{R}).\coqdoceol
\coqdocnoindent
\coqdockw{Proof}.\coqdoceol
\coqdocindent{1.00em}
\coqdoctac{intros} \coqdocvar{P} \coqdocvar{Q} \coqdocvar{R}. \coqdoctac{split}.\coqdoceol
\coqdocindent{1.00em}
\ensuremath{\times} \coqdoctac{intros} [\coqdocvar{HP} \ensuremath{|} [\coqdocvar{HQ} \coqdocvar{HR}]].\coqdoceol
\coqdocindent{2.00em}
** \coqdoctac{split}; \coqdoctac{left}; \coqdoctac{apply} \coqdocvar{HP}.\coqdoceol
\coqdocindent{2.00em}
** \coqdoctac{split}; \coqdoctac{right}. \coqdoctac{apply} \coqdocvar{HQ}. \coqdoctac{apply} \coqdocvar{HR}.\coqdoceol
\coqdocindent{1.00em}
\ensuremath{\times} \coqdoctac{intros} [[\coqdocvar{HP1} \ensuremath{|} \coqdocvar{HQ}] [\coqdocvar{HP2} \ensuremath{|} \coqdocvar{HR}]].\coqdoceol
\coqdocindent{2.00em}
** \coqdoctac{left}. \coqdoctac{apply} \coqdocvar{HP1}.\coqdoceol
\coqdocindent{2.00em}
** \coqdoctac{left}. \coqdoctac{apply} \coqdocvar{HP1}.\coqdoceol
\coqdocindent{2.00em}
** \coqdoctac{left}. \coqdoctac{apply} \coqdocvar{HP2}.\coqdoceol
\coqdocindent{2.00em}
** \coqdoctac{right}. \coqdoctac{split}. \coqdoctac{apply} \coqdocvar{HQ}. \coqdoctac{apply} \coqdocvar{HR}.\coqdoceol
\coqdocnoindent
\coqdockw{Qed}.\coqdoceol
\coqdocemptyline
\end{coqdoccode}
We are ignoring one of the hypotheses because it doesn't add anything to have two hypotheses which both simply state that P is true.

\begin{coqdoccode}
\coqdocemptyline
\coqdocnoindent
\coqdockw{Theorem} \coqdocvar{theorem\_1\_8\_iv'} : \coqdockw{\ensuremath{\forall}} (\coqdocvar{P} : \coqdockw{Prop}), \coqdocvar{P} \ensuremath{\land} \coqdocvar{P} \ensuremath{\leftrightarrow} \coqdocvar{P}.\coqdoceol
\coqdocnoindent
\coqdockw{Proof}.\coqdoceol
\coqdocindent{1.00em}
\coqdoctac{intros} \coqdocvar{P}. \coqdoctac{split}.\coqdoceol
\coqdocindent{1.00em}
\ensuremath{\times} \coqdoctac{intros} [\coqdocvar{HP} \coqdocvar{\_}].\coqdoceol
\coqdocindent{2.00em}
\coqdoctac{apply} \coqdocvar{HP}.\coqdoceol
\coqdocindent{1.00em}
\ensuremath{\times} \coqdoctac{intros} \coqdocvar{HP}. \coqdoctac{split}; \coqdoctac{apply} \coqdocvar{HP}.\coqdoceol
\coqdocnoindent
\coqdockw{Qed}.\coqdoceol
\end{coqdoccode}

  \coqlibrary{chapter 1 exercises 1 point 2}{Library }{chapter\_1\_exercises\_1\_point\_2}

\begin{coqdoccode}
\coqdocnoindent
\coqdockw{From} \coqdocvar{OUMBS} \coqdockw{Require} \coqdockw{Import} \coqdocvar{chapter\_1\_exercises\_1\_point\_1}.\coqdoceol
\coqdocemptyline
\coqdocnoindent
\coqdockw{Theorem} \coqdocvar{exercise\_1\_2\_a} : \coqdockw{\ensuremath{\forall}} (\coqdocvar{P} \coqdocvar{Q}: \coqdockw{Prop}), \~{}(\coqdocvar{P} \ensuremath{\lor} \coqdocvar{Q}) \ensuremath{\leftrightarrow} (\~{} \coqdocvar{P}) \ensuremath{\land} (\~{} \coqdocvar{Q}).\coqdoceol
\coqdocnoindent
\coqdockw{Proof}.\coqdoceol
\coqdocindent{1.00em}
\coqdoctac{intros} \coqdocvar{P} \coqdocvar{Q}.\coqdoceol
\coqdocindent{1.00em}
\coqdoctac{unfold} \coqdocvar{not}. \coqdoctac{split}.\coqdoceol
\coqdocindent{1.00em}
\ensuremath{\times} \coqdoctac{intros} \coqdocvar{HPQ}. \coqdoctac{split}.\coqdoceol
\coqdocindent{2.00em}
** \coqdoctac{intros} \coqdocvar{HP}. \coqdoctac{apply} \coqdocvar{or\_introl} \coqdockw{with} (\coqdocvar{B} := \coqdocvar{Q}) \coqdoctac{in} \coqdocvar{HP}. \coqdoctac{apply} \coqdocvar{HPQ} \coqdoctac{in} \coqdocvar{HP}. \coqdoctac{apply} \coqdocvar{HP}.\coqdoceol
\coqdocindent{2.00em}
** \coqdoctac{intros} \coqdocvar{HQ}. \coqdoctac{apply} \coqdocvar{or\_introl} \coqdockw{with} (\coqdocvar{B} := \coqdocvar{P}) \coqdoctac{in} \coqdocvar{HQ}. \coqdoctac{apply} \coqdocvar{theorem\_1\_8\_i} \coqdoctac{in} \coqdocvar{HQ}. \coqdoctac{apply} \coqdocvar{HPQ} \coqdoctac{in} \coqdocvar{HQ}. \coqdoctac{apply} \coqdocvar{HQ}.\coqdoceol
\coqdocindent{1.00em}
\ensuremath{\times} \coqdoctac{intros} [\coqdocvar{HNP} \coqdocvar{HNQ}]. \coqdoctac{intros} [\coqdocvar{HP} \ensuremath{|} \coqdocvar{HQ}]. \coqdoctac{apply} \coqdocvar{HNP} \coqdoctac{in} \coqdocvar{HP}. \coqdoctac{apply} \coqdocvar{HP}. \coqdoctac{apply} \coqdocvar{HNQ} \coqdoctac{in} \coqdocvar{HQ}. \coqdoctac{apply} \coqdocvar{HQ}.\coqdoceol
\coqdocnoindent
\coqdockw{Qed}.\coqdoceol
\coqdocemptyline
\coqdocnoindent
\coqdockw{Axiom} \coqdocvar{invalid\_middle} : \coqdockw{\ensuremath{\forall}} (\coqdocvar{P} : \coqdockw{Prop}), \coqdocvar{P} \ensuremath{\lor} \ensuremath{\lnot}\coqdocvar{P}.\coqdoceol
\coqdocemptyline
\coqdocnoindent
\coqdockw{Theorem} \coqdocvar{exercise\_1\_2\_b} : \coqdockw{\ensuremath{\forall}} (\coqdocvar{P} \coqdocvar{Q}: \coqdockw{Prop}), \~{}(\coqdocvar{P} \ensuremath{\land} \coqdocvar{Q}) \ensuremath{\leftrightarrow} (\~{} \coqdocvar{P}) \ensuremath{\lor} (\~{} \coqdocvar{Q}).\coqdoceol
\coqdocnoindent
\coqdockw{Proof}.\coqdoceol
\coqdocindent{1.00em}
\coqdoctac{intros} \coqdocvar{P} \coqdocvar{Q}.\coqdoceol
\coqdocindent{1.00em}
\coqdoctac{unfold} \coqdocvar{not}. \coqdoctac{split}.\coqdoceol
\coqdocindent{1.00em}
\ensuremath{\times} \coqdoctac{intros} \coqdocvar{HNPQ}. \coqdoctac{destruct} (\coqdocvar{invalid\_middle} \coqdocvar{P}) \coqdockw{as} [\coqdocvar{HP} \ensuremath{|} \coqdocvar{HNP}].\coqdoceol
\coqdocindent{2.00em}
** \coqdoctac{right}. \coqdoctac{intros} \coqdocvar{HQ}. \coqdoctac{assert} (\coqdocvar{H}: \coqdocvar{P} \ensuremath{\land} \coqdocvar{Q}). \{ \coqdoctac{split}. \coqdoctac{apply} \coqdocvar{HP}. \coqdoctac{apply} \coqdocvar{HQ}. \}\coqdoceol
\coqdocindent{3.50em}
\coqdoctac{apply} \coqdocvar{HNPQ} \coqdoctac{in} \coqdocvar{H}. \coqdoctac{apply} \coqdocvar{H}.\coqdoceol
\coqdocindent{2.00em}
** \coqdoctac{left}. \coqdoctac{apply} \coqdocvar{HNP}.\coqdoceol
\coqdocindent{1.00em}
\ensuremath{\times} \coqdoctac{intros} [\coqdocvar{HNP} \ensuremath{|} \coqdocvar{HNQ}].\coqdoceol
\coqdocindent{2.00em}
** \coqdoctac{intros} \coqdocvar{HPQ}. \coqdoctac{apply} \coqdocvar{proj1} \coqdoctac{in} \coqdocvar{HPQ}. \coqdoctac{apply} \coqdocvar{HNP} \coqdoctac{in} \coqdocvar{HPQ}. \coqdoctac{apply} \coqdocvar{HPQ}.\coqdoceol
\coqdocindent{2.00em}
** \coqdoctac{intros} \coqdocvar{HPQ}. \coqdoctac{apply} \coqdocvar{proj2} \coqdoctac{in} \coqdocvar{HPQ}. \coqdoctac{apply} \coqdocvar{HNQ} \coqdoctac{in} \coqdocvar{HPQ}. \coqdoctac{apply} \coqdocvar{HPQ}.\coqdoceol
\coqdocnoindent
\coqdockw{Qed}.\coqdoceol
\end{coqdoccode}

  \coqlibrary{chapter 1 exercises 1 point 3}{Library }{chapter\_1\_exercises\_1\_point\_3}

\begin{coqdoccode}
\coqdocnoindent
\coqdockw{From} \coqdocvar{OUMBS} \coqdockw{Require} \coqdockw{Import} \coqdocvar{chapter\_1\_exercises\_1\_point\_2}.\coqdoceol
\coqdocemptyline
\coqdocnoindent
\coqdockw{Theorem} \coqdocvar{exercise\_1\_3\_a} : \coqdockw{\ensuremath{\forall}} (\coqdocvar{P}: \coqdockw{Prop}), \~{}\~{}\coqdocvar{P} \ensuremath{\leftrightarrow} \coqdocvar{P}.\coqdoceol
\coqdocnoindent
\coqdockw{Proof}. \coqdoctac{split}.\coqdoceol
\coqdocnoindent
\coqdoctac{unfold} \coqdocvar{not}.\coqdoceol
\coqdocnoindent
\ensuremath{\times} \coqdoctac{intros} \coqdocvar{H}. \coqdoctac{destruct} (\coqdocvar{invalid\_middle} \coqdocvar{P}) \coqdockw{as} [\coqdocvar{HP} \ensuremath{|} \coqdocvar{HNP}]. \coqdoctac{apply} \coqdocvar{HP}. \coqdoctac{apply} \coqdocvar{H} \coqdoctac{in} \coqdocvar{HNP}. \coqdoctac{destruct} \coqdocvar{HNP}.\coqdoceol
\coqdocnoindent
\ensuremath{\times} \coqdoctac{intros} \coqdocvar{HP}. \coqdoctac{unfold} \coqdocvar{not}. \coqdoctac{intros} \coqdocvar{HNP}. \coqdoctac{apply} \coqdocvar{HNP} \coqdoctac{in} \coqdocvar{HP}. \coqdoctac{apply} \coqdocvar{HP}.\coqdoceol
\coqdocnoindent
\coqdockw{Qed}.\coqdoceol
\end{coqdoccode}

  \coqlibrary{chapter 1 exercises 1 point 4}{Library }{chapter\_1\_exercises\_1\_point\_4}

\begin{coqdoccode}
\coqdocnoindent
\coqdockw{From} \coqdocvar{OUMBS} \coqdockw{Require} \coqdockw{Import} \coqdocvar{chapter\_1\_exercises\_1\_point\_2}.\coqdoceol
\coqdocemptyline
\coqdocnoindent
\coqdockw{Theorem} \coqdocvar{exercise\_1\_4\_a} : \coqdockw{\ensuremath{\forall}} (\coqdocvar{P} \coqdocvar{Q}: \coqdockw{Prop}), (\coqdocvar{P} \ensuremath{\rightarrow} \coqdocvar{Q}) \ensuremath{\leftrightarrow} (\~{}\coqdocvar{Q} \ensuremath{\rightarrow} \ensuremath{\lnot}\coqdocvar{P}).\coqdoceol
\coqdocnoindent
\coqdockw{Proof}.\coqdoceol
\coqdocindent{1.00em}
\coqdoctac{intros} \coqdocvar{P} \coqdocvar{Q}. \coqdoctac{split}.\coqdoceol
\coqdocindent{1.00em}
\ensuremath{\times} \coqdoctac{intros} \coqdocvar{HPQ}. \coqdoctac{unfold} \coqdocvar{not}. \coqdoctac{intros} \coqdocvar{HNQ}. \coqdoctac{intros} \coqdocvar{HP}. \coqdoctac{apply} \coqdocvar{HPQ} \coqdoctac{in} \coqdocvar{HP}. \coqdoctac{apply} \coqdocvar{HNQ} \coqdoctac{in} \coqdocvar{HP}. \coqdoctac{apply} \coqdocvar{HP}.\coqdoceol
\coqdocindent{1.00em}
\ensuremath{\times} \coqdoctac{unfold} \coqdocvar{not}. \coqdoctac{intros} \coqdocvar{HNPNQ}. \coqdoctac{intros} \coqdocvar{HP}. \coqdoctac{destruct} (\coqdocvar{invalid\_middle} \coqdocvar{Q}).\coqdoceol
\coqdocindent{2.00em}
** \coqdoctac{apply} \coqdocvar{H}.\coqdoceol
\coqdocindent{2.00em}
** \coqdoctac{apply} \coqdocvar{HNPNQ} \coqdoctac{in} \coqdocvar{H}. \coqdoctac{destruct} \coqdocvar{H}. \coqdoctac{apply} \coqdocvar{HP}.\coqdoceol
\coqdocnoindent
\coqdockw{Qed}.\coqdoceol
\coqdocemptyline
\coqdocnoindent
\coqdockw{Theorem} \coqdocvar{exercise\_1\_4\_b} : \coqdockw{\ensuremath{\forall}} (\coqdocvar{P} \coqdocvar{Q}: \coqdockw{Prop}), (\coqdocvar{P} \ensuremath{\rightarrow} \coqdocvar{Q}) \ensuremath{\leftrightarrow} (\~{}\coqdocvar{P} \ensuremath{\lor} \coqdocvar{Q}).\coqdoceol
\coqdocnoindent
\coqdockw{Proof}.\coqdoceol
\coqdocindent{1.00em}
\coqdoctac{intros} \coqdocvar{P} \coqdocvar{Q}. \coqdoctac{split}.\coqdoceol
\coqdocindent{1.00em}
\ensuremath{\times} \coqdoctac{intros} \coqdocvar{HPQ}. \coqdoctac{destruct} (\coqdocvar{invalid\_middle} \coqdocvar{P}). \coqdoctac{apply} \coqdocvar{HPQ} \coqdoctac{in} \coqdocvar{H}. \coqdoctac{right}. \coqdoctac{apply} \coqdocvar{H}.\coqdoceol
\coqdocindent{2.00em}
\coqdoctac{left}. \coqdoctac{apply} \coqdocvar{H}.\coqdoceol
\coqdocindent{1.00em}
\ensuremath{\times} \coqdoctac{unfold} \coqdocvar{not}. \coqdoctac{intros} [\coqdocvar{HNP} \ensuremath{|} \coqdocvar{HQ}].\coqdoceol
\coqdocindent{2.00em}
** \coqdoctac{intros} \coqdocvar{HP}. \coqdoctac{apply} \coqdocvar{HNP} \coqdoctac{in} \coqdocvar{HP}. \coqdoctac{destruct} \coqdocvar{HP}.\coqdoceol
\coqdocindent{2.00em}
** \coqdoctac{intros} \coqdocvar{HP}. \coqdoctac{apply} \coqdocvar{HQ}.\coqdoceol
\coqdocnoindent
\coqdockw{Qed}.\coqdoceol
\coqdocemptyline
\coqdocnoindent
\coqdockw{Theorem} \coqdocvar{exercise\_1\_4\_c} : \coqdockw{\ensuremath{\forall}} (\coqdocvar{P} \coqdocvar{Q}: \coqdockw{Prop}), (\coqdocvar{P} \ensuremath{\rightarrow} \coqdocvar{Q}) \ensuremath{\leftrightarrow} \~{}(\coqdocvar{P} \ensuremath{\land} \ensuremath{\lnot}\coqdocvar{Q}).\coqdoceol
\coqdocnoindent
\coqdockw{Proof}.\coqdoceol
\coqdocindent{1.00em}
\coqdoctac{intros} \coqdocvar{P} \coqdocvar{Q}. \coqdoctac{split}.\coqdoceol
\coqdocindent{1.00em}
\ensuremath{\times} \coqdoctac{intros} \coqdocvar{HPQ}. \coqdoctac{unfold} \coqdocvar{not}. \coqdoctac{intros} [\coqdocvar{HP} \coqdocvar{HNQ}]. \coqdoctac{apply} \coqdocvar{HPQ} \coqdoctac{in} \coqdocvar{HP}. \coqdoctac{apply} \coqdocvar{HNQ} \coqdoctac{in} \coqdocvar{HP}. \coqdoctac{apply} \coqdocvar{HP}.\coqdoceol
\coqdocindent{1.00em}
\ensuremath{\times} \coqdoctac{unfold} \coqdocvar{not}. \coqdoctac{intros} \coqdocvar{HNPNQ} \coqdocvar{HP}. \coqdoctac{destruct} (\coqdocvar{invalid\_middle} \coqdocvar{Q}).\coqdoceol
\coqdocindent{2.00em}
** \coqdoctac{apply} \coqdocvar{H}.\coqdoceol
\coqdocindent{2.00em}
** \coqdoctac{assert} (\coqdocvar{H1}: \coqdocvar{P} \ensuremath{\land} \ensuremath{\lnot}\coqdocvar{Q}). \{ \coqdoctac{split}. \coqdoctac{apply} \coqdocvar{HP}. \coqdoctac{apply} \coqdocvar{H}. \} \coqdoctac{apply} \coqdocvar{HNPNQ} \coqdoctac{in} \coqdocvar{H1}. \coqdoctac{destruct} \coqdocvar{H1}.\coqdoceol
\coqdocnoindent
\coqdockw{Qed}.\coqdoceol
\coqdocemptyline
\coqdocnoindent
\coqdockw{Theorem} \coqdocvar{exercise\_1\_4\_d} : \coqdockw{\ensuremath{\forall}} (\coqdocvar{P} \coqdocvar{Q}: \coqdockw{Prop}), (\coqdocvar{P} \ensuremath{\land} (\coqdocvar{P} \ensuremath{\rightarrow} \coqdocvar{Q})) \ensuremath{\rightarrow} \coqdocvar{Q}.\coqdoceol
\coqdocnoindent
\coqdockw{Proof}.\coqdoceol
\coqdocindent{1.00em}
\coqdoctac{intros} \coqdocvar{P} \coqdocvar{Q} [\coqdocvar{HP} \coqdocvar{HPQ}].\coqdoceol
\coqdocindent{1.00em}
\coqdoctac{apply} \coqdocvar{HPQ} \coqdoctac{in} \coqdocvar{HP}. \coqdoctac{apply} \coqdocvar{HP}.\coqdoceol
\coqdocnoindent
\coqdockw{Qed}.\coqdoceol
\coqdocemptyline
\coqdocnoindent
\coqdockw{Theorem} \coqdocvar{exercise\_1\_4\_e} : \coqdockw{\ensuremath{\forall}} (\coqdocvar{P} \coqdocvar{Q}: \coqdockw{Prop}), ((\coqdocvar{P} \ensuremath{\lor} \coqdocvar{Q}) \ensuremath{\land} \ensuremath{\lnot}\coqdocvar{P}) \ensuremath{\rightarrow} \coqdocvar{Q}.\coqdoceol
\coqdocnoindent
\coqdockw{Proof}.\coqdoceol
\coqdocindent{1.00em}
\coqdoctac{intros} \coqdocvar{P} \coqdocvar{Q} [[\coqdocvar{HP} \ensuremath{|} \coqdocvar{HQ}] \coqdocvar{HNP}].\coqdoceol
\coqdocindent{1.00em}
\ensuremath{\times} \coqdoctac{apply} \coqdocvar{HNP} \coqdoctac{in} \coqdocvar{HP}. \coqdoctac{destruct} \coqdocvar{HP}.\coqdoceol
\coqdocindent{1.00em}
\ensuremath{\times} \coqdoctac{apply} \coqdocvar{HQ}.\coqdoceol
\coqdocnoindent
\coqdockw{Qed}.\coqdoceol
\end{coqdoccode}

  \coqlibrary{chapter 1 exercises 1 point 5}{Library }{chapter\_1\_exercises\_1\_point\_5}

\begin{coqdoccode}
\coqdocnoindent
\coqdockw{From} \coqdocvar{OUMBS} \coqdockw{Require} \coqdockw{Import} \coqdocvar{chapter\_1\_exercises\_1\_point\_1}.\coqdoceol
\coqdocemptyline
\coqdocnoindent
\coqdockw{Theorem} \coqdocvar{exercise\_1\_5\_a} : \coqdockw{\ensuremath{\forall}} (\coqdocvar{P} \coqdocvar{Q} \coqdocvar{R} : \coqdockw{Prop}), ((\coqdocvar{P} \ensuremath{\rightarrow} \coqdocvar{Q}) \ensuremath{\land}  (\coqdocvar{Q} \ensuremath{\rightarrow} \coqdocvar{R})) \ensuremath{\rightarrow} (\coqdocvar{P} \ensuremath{\rightarrow} \coqdocvar{R}).\coqdoceol
\coqdocnoindent
\coqdockw{Proof}.\coqdoceol
\coqdocindent{1.00em}
\coqdoctac{intros} \coqdocvar{P} \coqdocvar{Q} \coqdocvar{R} [\coqdocvar{HPQ} \coqdocvar{HQR}] \coqdocvar{H}.\coqdoceol
\coqdocindent{1.00em}
\coqdoctac{apply} \coqdocvar{HPQ} \coqdoctac{in} \coqdocvar{H}. \coqdoctac{apply} \coqdocvar{HQR} \coqdoctac{in} \coqdocvar{H}. \coqdoctac{apply} \coqdocvar{H}.\coqdoceol
\coqdocnoindent
\coqdockw{Qed}.\coqdoceol
\coqdocemptyline
\coqdocnoindent
\coqdockw{Theorem} \coqdocvar{exercise\_1\_5\_b} : \coqdockw{\ensuremath{\forall}} (\coqdocvar{P} \coqdocvar{Q} \coqdocvar{R} : \coqdockw{Prop}), ((\coqdocvar{P} \ensuremath{\rightarrow} \coqdocvar{Q}) \ensuremath{\land}  (\coqdocvar{R} \ensuremath{\rightarrow} \coqdocvar{Q})) \ensuremath{\leftrightarrow} ((\coqdocvar{P} \ensuremath{\lor} \coqdocvar{R}) \ensuremath{\rightarrow} \coqdocvar{Q}).\coqdoceol
\coqdocnoindent
\coqdockw{Proof}.\coqdoceol
\coqdocindent{1.00em}
\coqdoctac{intros} \coqdocvar{P} \coqdocvar{Q} \coqdocvar{R}. \coqdoctac{split}.\coqdoceol
\coqdocindent{1.00em}
\ensuremath{\times} \coqdoctac{intros} [\coqdocvar{HPQ} \coqdocvar{HQR}] [\coqdocvar{HP} \ensuremath{|} \coqdocvar{HR}].\coqdoceol
\coqdocindent{2.00em}
** \coqdoctac{apply} \coqdocvar{HPQ} \coqdoctac{in} \coqdocvar{HP}. \coqdoctac{apply} \coqdocvar{HP}.\coqdoceol
\coqdocindent{2.00em}
** \coqdoctac{apply} \coqdocvar{HQR} \coqdoctac{in} \coqdocvar{HR}. \coqdoctac{apply} \coqdocvar{HR}.\coqdoceol
\coqdocindent{1.00em}
\ensuremath{\times} \coqdoctac{intros} \coqdocvar{HPRQ}. \coqdoctac{split}.\coqdoceol
\coqdocindent{2.00em}
** \coqdoctac{intros} \coqdocvar{HP}. \coqdoctac{apply} \coqdocvar{or\_introl} \coqdockw{with} (\coqdocvar{B} := \coqdocvar{R}) \coqdoctac{in} \coqdocvar{HP}. \coqdoctac{apply} \coqdocvar{HPRQ}. \coqdoctac{apply} \coqdocvar{HP}.\coqdoceol
\coqdocindent{2.00em}
** \coqdoctac{intros} \coqdocvar{HR}. \coqdoctac{apply} \coqdocvar{or\_introl} \coqdockw{with} (\coqdocvar{B} := \coqdocvar{P}) \coqdoctac{in} \coqdocvar{HR}. \coqdoctac{apply} \coqdocvar{theorem\_1\_8\_i} \coqdoctac{in} \coqdocvar{HR}. \coqdoctac{apply} \coqdocvar{HPRQ}. \coqdoctac{apply} \coqdocvar{HR}.\coqdoceol
\coqdocnoindent
\coqdockw{Qed}.\coqdoceol
\coqdocemptyline
\coqdocnoindent
\coqdockw{Theorem} \coqdocvar{exercise\_1\_5\_c} : \coqdockw{\ensuremath{\forall}} (\coqdocvar{P} \coqdocvar{Q} \coqdocvar{R} : \coqdockw{Prop}), ((\coqdocvar{P} \ensuremath{\rightarrow} \coqdocvar{Q}) \ensuremath{\land}  (\coqdocvar{P} \ensuremath{\rightarrow} \coqdocvar{R})) \ensuremath{\leftrightarrow} (\coqdocvar{P} \ensuremath{\rightarrow} (\coqdocvar{Q} \ensuremath{\land} \coqdocvar{R})).\coqdoceol
\coqdocnoindent
\coqdockw{Proof}.\coqdoceol
\coqdocindent{1.00em}
\coqdoctac{intros} \coqdocvar{P} \coqdocvar{Q} \coqdocvar{R}. \coqdoctac{split}.\coqdoceol
\coqdocindent{1.00em}
\ensuremath{\times} \coqdoctac{intros} [\coqdocvar{HPQ} \coqdocvar{HPR}] \coqdocvar{HP}. \coqdoctac{split}.\coqdoceol
\coqdocindent{2.00em}
** \coqdoctac{apply} \coqdocvar{HPQ}. \coqdoctac{apply} \coqdocvar{HP}.\coqdoceol
\coqdocindent{2.00em}
** \coqdoctac{apply} \coqdocvar{HPR}. \coqdoctac{apply} \coqdocvar{HP}.\coqdoceol
\coqdocindent{1.00em}
\ensuremath{\times} \coqdoctac{intros} \coqdocvar{HPQR}. \coqdoctac{split}; \coqdoctac{intros} \coqdocvar{HP}; \coqdoctac{apply} \coqdocvar{HPQR}; \coqdoctac{apply} \coqdocvar{HP}.\coqdoceol
\coqdocnoindent
\coqdockw{Qed}.\coqdoceol
\end{coqdoccode}

  \coqlibrary{chapter 1 exercises 1 point 6}{Library }{chapter\_1\_exercises\_1\_point\_6}

\begin{coqdoccode}
\coqdocnoindent
\coqdockw{Theorem} \coqdocvar{exercise\_1\_6\_a} : \coqdockw{\ensuremath{\forall}} (\coqdocvar{P} \coqdocvar{Q} \coqdocvar{R} : \coqdockw{Prop}), (\coqdocvar{Q} \ensuremath{\leftrightarrow} \coqdocvar{R}) \ensuremath{\rightarrow} (\coqdocvar{P} \ensuremath{\lor} \coqdocvar{Q} \ensuremath{\leftrightarrow} \coqdocvar{P} \ensuremath{\lor} \coqdocvar{R}).\coqdoceol
\coqdocnoindent
\coqdockw{Proof}.\coqdoceol
\coqdocindent{1.00em}
\coqdoctac{intros} \coqdocvar{P} \coqdocvar{Q} \coqdocvar{R} \coqdocvar{HQR}. \coqdoctac{split}.\coqdoceol
\coqdocindent{1.00em}
\ensuremath{\times} \coqdoctac{intros} [\coqdocvar{HP} \ensuremath{|} \coqdocvar{HQ}].\coqdoceol
\coqdocindent{2.00em}
** \coqdoctac{left}. \coqdoctac{apply} \coqdocvar{HP}.\coqdoceol
\coqdocindent{2.00em}
** \coqdoctac{right}. \coqdoctac{apply} \coqdocvar{HQR}. \coqdoctac{apply} \coqdocvar{HQ}.\coqdoceol
\coqdocindent{1.00em}
\ensuremath{\times} \coqdoctac{intros} [\coqdocvar{HP} \ensuremath{|} \coqdocvar{HR}].\coqdoceol
\coqdocindent{2.00em}
** \coqdoctac{left}. \coqdoctac{apply} \coqdocvar{HP}.\coqdoceol
\coqdocindent{2.00em}
** \coqdoctac{right}. \coqdoctac{apply} \coqdocvar{HQR}. \coqdoctac{apply} \coqdocvar{HR}.\coqdoceol
\coqdocnoindent
\coqdockw{Qed}.\coqdoceol
\coqdocemptyline
\coqdocnoindent
\coqdockw{Theorem} \coqdocvar{exercise\_1\_6\_b} : \coqdockw{\ensuremath{\forall}} (\coqdocvar{P} \coqdocvar{Q} \coqdocvar{R} : \coqdockw{Prop}), (\coqdocvar{Q} \ensuremath{\leftrightarrow} \coqdocvar{R}) \ensuremath{\rightarrow} (\coqdocvar{P} \ensuremath{\land} \coqdocvar{Q} \ensuremath{\leftrightarrow} \coqdocvar{P} \ensuremath{\land} \coqdocvar{R}).\coqdoceol
\coqdocnoindent
\coqdockw{Proof}.\coqdoceol
\coqdocindent{1.00em}
\coqdoctac{intros} \coqdocvar{P} \coqdocvar{Q} \coqdocvar{R} \coqdocvar{HQR}. \coqdoctac{split}.\coqdoceol
\coqdocindent{1.00em}
\ensuremath{\times} \coqdoctac{intros} [\coqdocvar{HP} \coqdocvar{HQ}]. \coqdoctac{split}.\coqdoceol
\coqdocindent{2.00em}
** \coqdoctac{apply} \coqdocvar{HP}.\coqdoceol
\coqdocindent{2.00em}
** \coqdoctac{apply} \coqdocvar{HQR}. \coqdoctac{apply} \coqdocvar{HQ}.\coqdoceol
\coqdocindent{1.00em}
\ensuremath{\times} \coqdoctac{intros} [\coqdocvar{HP} \coqdocvar{HR}]. \coqdoctac{split}.\coqdoceol
\coqdocindent{2.00em}
** \coqdoctac{apply} \coqdocvar{HP}.\coqdoceol
\coqdocindent{2.00em}
** \coqdoctac{apply} \coqdocvar{HQR}. \coqdoctac{apply} \coqdocvar{HR}.\coqdoceol
\coqdocnoindent
\coqdockw{Qed}.\coqdoceol
\coqdocemptyline
\coqdocnoindent
\coqdockw{Theorem} \coqdocvar{exercise\_1\_6\_c} : \coqdockw{\ensuremath{\forall}} (\coqdocvar{P} \coqdocvar{Q} \coqdocvar{R} : \coqdockw{Prop}), (\coqdocvar{Q} \ensuremath{\leftrightarrow} \coqdocvar{R}) \ensuremath{\rightarrow} ((\coqdocvar{P} \ensuremath{\rightarrow} \coqdocvar{Q}) \ensuremath{\leftrightarrow} (\coqdocvar{P} \ensuremath{\rightarrow} \coqdocvar{R})).\coqdoceol
\coqdocnoindent
\coqdockw{Proof}.\coqdoceol
\coqdocindent{1.00em}
\coqdoctac{intros} \coqdocvar{P} \coqdocvar{Q} \coqdocvar{R} \coqdocvar{HQR}. \coqdoctac{split}.\coqdoceol
\coqdocindent{1.00em}
\ensuremath{\times} \coqdoctac{intros} \coqdocvar{HPQ} \coqdocvar{HP}. \coqdoctac{apply} \coqdocvar{HQR}. \coqdoctac{apply} \coqdocvar{HPQ}. \coqdoctac{apply} \coqdocvar{HP}.\coqdoceol
\coqdocindent{1.00em}
\ensuremath{\times} \coqdoctac{intros} \coqdocvar{HPR} \coqdocvar{HP}. \coqdoctac{apply} \coqdocvar{HQR}. \coqdoctac{apply} \coqdocvar{HPR}. \coqdoctac{apply} \coqdocvar{HP}.\coqdoceol
\coqdocnoindent
\coqdockw{Qed}.\coqdoceol
\end{coqdoccode}


 \subsection*{Exercises 1.2}
		 
		 \begin{enumerate}
		 	\item[1.] \begin{enumerate}
		 		\item[a)] $A \cup C \subseteq B \cup D$
		 		
		 		\begin{proof}
		 		First, suppose that $A \subseteq B$ and $C \subseteq D$, that is by definition that $x \in~A~\Rightarrow~x~\in~B$ and $x \in C \Rightarrow x \in D$. Then, by exercise 1.1.7(a) we can say that $x \in A~\quad\lor \quad~x~\in~C \Rightarrow x \in B \lor x \in D$ and by definition 1.13 we got $x \in A \cup C \Rightarrow x \in B \cup D$. Finally, we can conclude $A \cup C \subseteq B \cup D$ by definition 1.10.  
		 		\end{proof}
		 		
		 		\item[b)] Analogous to the prevoius exercise, using execersice 1.1.7(b).
		 	\end{enumerate}
		 	
		 	\item[2.] \begin{enumerate}
		 		\item[a)] $A \cup C = B \cup D$
		 		
		 		\begin{proof}
		 			First, suppose that $A = B$ and $C = D$. We have to prove just that $B \cup D \subseteq A \cup C$ since $A \cup C \subseteq B \cup D$ was proved in the preceding exercise, having in mind that the hypothesis of the excerise is met, because we are supposing $A = B$ and $C = D$, so, we got by the axiom of extent $x \in A \Rightarrow x \in B$, $x \in B \Rightarrow x \in A$ and $x \in C \Rightarrow x \in D$, $x \in D \Rightarrow x \in C$, that is, $A \subseteq B$, $B \subseteq A$ and $C \subseteq D$, $D \subseteq C$. Then, because of $x \in B \Rightarrow x \in A$ and $x \in D \Rightarrow x \in C$, we have 
		 			\[
		 			x \in B \lor x \in D \Rightarrow x \in A \lor x \in C
		 			\]
		 			by exercise 1.1.7(a) and by definition 1.13 we got $x \in B \cup D \Rightarrow x \in A \cup C$. Finally, we can say that $B \cup D \subseteq A \cup C$ and $A \cup C \subseteq B \cup D$ and by theorem 1.11(iv) we conclude that $A \cup C = B \cup D$.
		 		\end{proof}
		 		
		 		\item[b)] Analogous to the prevoius exercise, using execersice 1.1.7(b).
		 	\end{enumerate}
		 	
		 	\item[3.] \begin{proof}
		 		Let's suppose $A \subseteq B$. Then we have $x \in A \Rightarrow x \in B$ by the definition 1.10 and since $(P \Rightarrow Q) \Leftrightarrow (\neg Q \Rightarrow \neg P)$ by exercise 1.1.4(a), we got $x \notin B \Rightarrow x \notin A$, that is, $x \in B' \Rightarrow x \in A'$ by definition 1.19 and finally, we can conclude by definition 1.10 that $B' \subseteq A'$.
		 	\end{proof}
		 	
		 	\item[4.] \begin{proof}
		 		Since $A = B$, we have that $A \subseteq B$ and $B \subseteq A$ as a consequence of the axiom of extent and then, we can apply the preceding exercise and the theorem 1.11 (iv) to conclude $A' = B'$.
		 	\end{proof}
		 	
		 	\item[5.] \begin{proof}
		 		Suppose $A = B$ and $B \subseteq C$. We know by the axiom of extent that $x \in A \Rightarrow x \in B$ (i) and $x \in B \Rightarrow x \in A$. Also, we have $x \in B \Rightarrow x \in C$ and, because of exercise 1.1.5(a) we can say that $x \in A \Rightarrow x \in C$ and then we can conclude $A \subseteq C$ by definition 1.10.
		 	\end{proof}
		 	
		 	\item[6.] \begin{proof}
		 	 Suppose $A \subset B$ and $B \subset C$. We want prove that $x \in A \Rightarrow x \in C$ and $A \neq C$. By the hypothesis we have $x \in A \Rightarrow x \in B$ and $x \in B \Rightarrow x \in C$ (i), with $A \neq B \neq C$ and then, applying exercise 1.5(a) we can say that $x \in A \Rightarrow x \in C$. Now, let's suppose $A = C$. Since $C = A \subset B$, we have by definition 1.10 that $x \in C \Rightarrow x \in B$ and if we use that fact and (i), applying the axiom of extent we got that $B = C$ (absurd). Finally, we can conclude $x \in A \Rightarrow x \in C$ and $A \neq C$, that is, $A \subset C$.\end{proof}
		 \end{enumerate}
		
		\item[7.] \begin{enumerate}
			\item[iv)] \begin{proof}Since $A \subseteq B$ and $B \subseteq A$ is by definition 1.10 equivalent to say $x \in A \Rightarrow x \in B$ and $x \in B \Rightarrow x \in A$, we have by de axiom of extent that $A = B$. \end{proof}
			\item[v)] Just use the definition and apply exercise 1.1.5(a).
		\end{enumerate}
		
		\item[8.] Suppose that $S \in S$. By definition of $S$ we will have that $S \notin S$, which is a contradiction. Then we can conclude that $S$ is not an element.
		
		\item[9.] No, because A2 just permit us to form the class of all elements that satisfy a property $P(x)$, not a class of classes that satisfy the property.
		
		\item[10.] It cannot be produced because A2 just allow to construct classes which its elements satisfy a condition, not to form a class of classes that satisfy some property.
	\end{enumerate}

\subsection*{Exercises 1.3}
	
\end{document}
	
	
	
\end{document}
